\tableofcontents

\section{Career development}

\noindent
\textbf{How would you describe your level of experience as a professional software engineer?
What are your strengths as a software engineer?}

I have worked as a software engineer for more than ten years, and I consider myself a senior engineer.
I started at university by writing C for embedded devices, kernel modules for FPGA hardware interfaces and
a FAT file system for our monolithic kernel project. Then I wrote an IoT platform that talked to devices over
LoRaWAN, Zigbee and nRF, and that moved me to backend development in Go.

From 2019 to 2025 I worked at \href{https://snapp.ir}{Snapp!}, Iran's largest ride-hailing platform with more than 50 million users.
I started as a software engineer in the shared services team and became a senior cloud engineer and then lead engineer
for platform architecture. I led the design of the company-wide cloud platform on Kubernetes, wrote Kubernetes operators
(CRDs) for ArgoCD authentication and authorization, and designed a microservice architecture on NATS that handled more than
300k messages per second. Since January 2025 I have been a senior software engineer at AveeHealth, a Canadian digital-health
startup, where I build the product in Python and run the platform on AWS with EKS and ECS. In 2026 I moved to Barcelona.

My main strengths are:

\begin{itemize}
  \item \textbf{Academic knowledge together with production experience.} I have a B.Sc., M.Sc. and Ph.D. from Amirkabir University
    of Technology, and I have taught Internet Engineering there for seven semesters. I try to understand a problem from first principles
    before I fix it, because computers are logical and the root cause is always there to find.
  \item \textbf{Operating what I build.} I was on call at Snapp! and I am on call at AveeHealth. I have detected issues, led incident
    calls and made decisions under pressure.
  \item \textbf{Adaptability.} At Snapp! I introduced NATS as the company's messaging system and ran workshops so other teams could
    use it. When the company needed machine learning pipelines, I built them with KServe, Knative and Airflow, and a FastAPI serving
    layer so the data science team could deploy their own models.
  \item \textbf{Domain knowledge.} I always try to learn the domain of the product. When I worked on software for the oil industry,
    I learned how pumpjacks and oil wells work, because good software needs domain knowledge.
\end{itemize}

\noindent
\textbf{What are your favourite software development environments, languages, and tools? Elaborate on your choices.}

I use \href{https://neovim.io/}{Neovim} every day, because it is fast and I can customise it fully. I keep my editor and my whole
system configuration in a dotfiles framework that I maintain, and I test it with fresh installations of Arch Linux.

Go is my main language. When I first saw it I thought of it as C for high-level services: simple, fast, with concurrency built in.
Python is my second language and most of my recent product work is in it (FastAPI, Django, Pydantic AI). I use C when I need to work
close to the hardware or the kernel, and I have also learned Rust. Over the years I have had to write PHP, Scala, Java, Node.js and even Perl.

My daily tools are Docker and Podman, Git with GitHub and GitLab CI, kubectl, Helm, \href{https://argo-cd.readthedocs.io/}{ArgoCD},
stern and krew, Ansible and Terraform, tmux, and Arch Linux on my machines with Ubuntu on servers. I prefer command-line tools, containers
instead of installing dependencies on the host, and projects that are truly open source. I also maintain packages on the Arch User Repository.

\noindent
\textbf{How extensive is your experience of C++ software engineering? Outline the applications that you have led in C++,
and your takeaways from that experience.}

My low-level experience is mostly C, with some C++. I took part in the \href{http://eudyptula-challenge.org/}{Eudyptula Challenge},
implemented a FAT file system in C, and in the university IoT working group I designed a reliable layer-2 transport protocol
(stop-and-wait ARQ with retries) on top of raw nRF radio modules. In C++ I worked on a Qt touch dashboard for our home automation
system on a Raspberry Pi, firmware for LoRaWAN farm sensors, and changes to our private forks of the
\href{https://github.com/valhalla/valhalla}{Valhalla} and \href{https://github.com/Project-OSRM/osrm-backend}{OSRM} routing engines at Snapp!.

The main lessons were about memory and threads: we used \href{https://valgrind.org/}{Valgrind} regularly to find leaks, and we had to deal with
binary protocol headers and endianness, cross-compilation between Intel and ARM, and memory allocation on devices without an operating system.
For the last years my main languages have been Go and Python, so I would describe my C++ as solid but not current.

\noindent
\textbf{Describe top Linux projects (incl. containers) you have worked on (purpose, market, etc.).
What are your contributions (design, implementation, testing, documentation, CI/CD, etc)?}

\begin{itemize}
  \item \textbf{Snapp! cloud platform.} I led the design and implementation of the company-wide platform on Kubernetes. I wrote
    Kubernetes operators (CRDs) that automated ArgoCD authentication and authorization, deployed Kafka with the
    \href{https://strimzi.io/}{Strimzi} operator and RabbitMQ and Redis through their operators and custom resources, and set up monitoring
    and alerting with OpenTelemetry. For Baly, Snapp!'s branch in Iraq, we ran services on AWS with EKS and Terraform.
  \item \textbf{Surge.} The real-time pricing service at Snapp!. I led it from the beginning: we divided cities into districts, one part consumed
    driver-location and price-request events and did 15-minute windowing before publishing to our NATS cluster, and a second part read the batches
    and calculated the price change. It increased accepted rides by 5\%. The QA team tested it, and we provided the Kubernetes environments they needed.
  \item \textbf{ML platform.} I built ML pipelines with KServe and Knative on Kubernetes, a Python/FastAPI serving layer so the data science team could
    deploy models from notebook to production, an \href{https://ollama.com/}{Ollama} cluster with LiteLLM in front of it to share load between local
    and third-party LLMs, and Airflow in place of legacy cron jobs.
  \item \textbf{AveeHealth.} I run the platform on AWS with managed EKS clusters and ECS, and deploy the services in Docker containers.
  \item \textbf{Device security agent (2025).} An on-device agent in Go for a US security startup, with a device-to-backend channel over a message broker,
    deployed with Docker and Kubernetes.
  \item \textbf{Open source.} I contributed to \href{https://github.com/nats-io/}{NATS} and \href{https://github.com/emqx}{EMQX}, mostly to their
    Kubernetes operators and Helm charts.
\end{itemize}

In these projects I did the design and most of the implementation, and I designed our CI/CD pipelines on GitLab and GitHub.

\section{Experience}

\noindent
\textbf{Please explain the differences between running an application natively, in a container, or VM.}

Natively, the application runs directly on the host and needs every library and dependency installed there, so different applications can conflict.
A container packages the application with its user space and isolates it with Linux kernel features (namespaces, cgroups and a layered file system),
but it shares the host kernel. A VM emulates a whole machine with its own kernel, so it uses more resources but isolates at a lower level and can run
a different operating system. When you need to change kernel settings or run another operating system, a container is not enough.

On Kubernetes this difference shows up every day: containers make it cheap to schedule, restart and scale workloads, but noisy neighbours, kernel
versions and resource limits are shared concerns. For example, I debugged an Ollama upgrade on Kubernetes where pods started normally but connections
failed, and the cause was the cluster kernel together with Go's MPTCP behaviour. The fix was on the node kernel side, not in the container.

\noindent
\textbf{Outline your thoughts on open source software.
What is important to get right when you are working on open source projects?}

Open source is the most open way to build software: anyone can contribute, whatever their nationality or background, and projects grow faster because
of that. I started an open source organisation at each company I worked for and released the tools that could be useful to others on GitHub. I maintain
packages on the AUR and I have contributed to NATS and EMQX.

The most important things are documentation, so people can use and extend the project, and maintenance: fix bugs and answer issues, or archive the
project clearly. Scope also matters. When people ask for features outside the roadmap, the maintainers should say so clearly instead of leaving the
issue open.

\noindent
\textbf{Outline your thoughts on performance in software engineering. How do you ensure that your team's code is performant?
Can you describe hard learning in this regard?}

At Snapp! we had millions of requests per day, and any commit could add a memory leak or increase response time. We ran a load test with
\href{https://github.com/grafana/k6}{k6} for each release, and after the release we watched Prometheus and Grafana dashboards with response-time
histograms (p50, p90, p99). For more detail we used tracing and continuous profiling in staging.

The hard lesson came on a smaller B2B project: the team thought these tools were unnecessary, then queries in production took more than 30 seconds and
AWS closed the connections, and we had no data to find which queries were slow. Since then I add metrics and tracing from the first version, even for
small services. In Go specifically, the most common performance problems I have seen are unbounded concurrency and missing backpressure, so I limit
concurrency explicitly and use context for cancellation and timeouts.

\noindent
\textbf{Outline your thoughts on documentation in software projects. What practices do you think software developers
should follow in this regard? What are examples of the best open source documentation you have worked with?}

Documentation should live next to the code so it is updated in the same review. Outdated documentation is worse than none, because people follow it
instead of asking. I prefer comments in the code, a README that explains how to run and test the project, contribution guidelines, and a short
high-level design. For bigger projects, an RFC process lets the team track architectural and protocol changes and reference them from the code.
Schemas should be part of the protocol where possible, for example with Protobuf or Avro, so the schema is the documentation.

The best documentation I have used is \href{https://docs.djangoproject.com/}{Django}'s: every feature has at least one working example that is easy to read.

\section{Education}

\noindent
\textbf{How did you rank in your final year of high school in mathematics? Were you a top student? On what basis would you say that?}

I was in the top 10\% of my school in mathematics and computer programming, based on the final-year report our high school gave us.

\noindent
\textbf{How did you rank in your final year of high school, in your home language? Were you a top student? On what basis would you say that?}

My home language is Persian, and I was not a top student in it.

\noindent
\textbf{Please state your high school graduation results or university entrance results, and explain the grading system used.}

Iran has one national university entrance exam (Konkur). I ranked 646 out of about 100,000 students in my year, which got me into Amirkabir University
of Technology, one of the top technical universities in Iran.

\noindent
\textbf{Can you make a case that you are in the top 5\% in your academic year, or top 1\%, or even higher?}

The entrance exam rank, 646 of about 100,000, is within the top 1\% nationally. At university I was first in my class in both my B.Sc. (18.94 out of 20)
and my M.Sc. (18.77 out of 20); Iranian universities grade from 0 to 20, with 20 the best. I then completed a Ph.D. in Computer Networks (17.34 out of 20).
References are available on request.

\noindent
\textbf{What sort of high school student were you? Outside of class, what were your interests and hobbies? What would your high school peers remember you for?}

I was a bookworm and spent most of my free time in the school library. I was also interested in computer networks, and my friends remember me for
setting up their computers on a local network so we could play Counter-Strike together.

\noindent
\textbf{Which university and degree did you choose? What other universities did you consider, and why did you select that one?}

I fell in love with computers, and especially networks, in high school, so I chose computer engineering. I chose Amirkabir University of Technology because
it is one of the most respected technical universities in Iran, known as the mother of the technical universities.

\noindent
\textbf{Overall, what was your degree result and how did that reflect on your ability?}

I was first in my class in my B.Sc. and M.Sc. I think this came from real passion for operating systems and networks. What I learned there still helps
me every day: I try to understand the problem and reason about it logically before I act, and that is how I usually find the root cause.

\noindent
\textbf{During all of your education years, from high school to university, can you describe any achievements that were truly exceptional?}

\begin{itemize}
  \item Ranked 4th in the first round and 10th in the second round of the National Olympiad of Computer Networks during my B.Sc.
  \item Passed two rounds of the National Olympiad in Informatics in high school.
  \item Our team was the first from my university to take part in the 39th ACM International Collegiate Programming Contest (West Asia region).
  \item Ranked 2nd in the national hardware-software co-design competition, implementing a genetic algorithm in hardware and software.
\end{itemize}

\noindent
\textbf{What leadership roles did you take on during your education?
Did you conceive of, and drive to completion, any initiatives outside of your required classwork?}

I led the technical side of events in our Computer Engineering faculty, including online and onsite programming contests and Linux and open source
festivals, and an open source project that used \href{https://freeradius.org/}{FreeRADIUS} to manage internet usage for students and staff.
I was a teaching assistant for at least one course every semester of my B.Sc., and after my M.Sc. the university invited me to teach. I taught Internet
Engineering for seven semesters and Principles of Computer and Programming for three.

\section{Context}

\noindent
\textbf{Outline your thoughts on the mission of Canonical. What is it about the company's purpose and goals which is most appealing to you?
What is it that you think is risky or unappealing?}

Canonical's mission, as I see it, is to put open source in the hands of everyone, from a person installing Linux for the first time to an enterprise
running Kubernetes in production. Ubuntu was my own first step into Linux as a child, because it was easy to use and well documented. The same idea now
applies to infrastructure: Charmed operators and Canonical Kubernetes give people a production platform that is free, secure and supported.

The risk is the business model: many companies use open source without paying for it. I think it works when the product is good enough that operating it
yourself costs more than buying support, and that is exactly where operators and lifecycle management add value.

\noindent
\textbf{What would you like to achieve in career and skills development?}

I want to work on Kubernetes itself rather than only on top of it. I have built operators and run platforms for my own companies, and building the
operators and add-ons that other companies depend on is the next step. I also want to have a bigger effect on the open source world, and working in
the open at Canonical, with a community and upstream projects, is the best place I know to do that.
